\documentclass[11pt,a4paper]{article}

% ---------------------------------------------------------------------------
% PACKAGES
% ---------------------------------------------------------------------------
\usepackage[utf8]{inputenc}
\usepackage[T1]{fontenc}
\usepackage[french]{babel}
\usepackage{amsmath}
\usepackage{amssymb}
\usepackage{lmodern}
\usepackage{geometry}
\usepackage{xcolor}
\usepackage{booktabs}
\usepackage{longtable}
\usepackage{array}
\usepackage{tabularx}
\usepackage{graphicx}
\usepackage{fancyhdr}
\usepackage{titlesec}
\usepackage{caption}
\usepackage{enumitem}
\usepackage[hidelinks]{hyperref}
\usepackage{listings}

% ---------------------------------------------------------------------------
% GEOMETRY & PAGE
% ---------------------------------------------------------------------------
\geometry{margin=2.2cm,headheight=22pt}
\setlength{\parindent}{0pt}
\setlength{\parskip}{6pt}

\definecolor{vocblue}{RGB}{0,102,153}
\definecolor{vocdark}{RGB}{20,40,60}
\definecolor{voclight}{RGB}{235,242,247}
\definecolor{vocgreen}{RGB}{0,140,90}
\definecolor{vocred}{RGB}{180,40,40}
\definecolor{vocorange}{RGB}{210,120,0}

\hypersetup{
  colorlinks=true,
  linkcolor=vocblue,
  urlcolor=vocblue,
  citecolor=vocblue,
  pdftitle={Rapport de stage -- VOC Platform},
  pdfauthor={TBINI Mustapha Amin}
}

% ---------------------------------------------------------------------------
% HEADERS & FOOTERS
% ---------------------------------------------------------------------------
\pagestyle{fancy}
\fancyhf{}
\fancyhead[L]{\small\textcolor{vocblue}{\textbf{VOC Platform -- Vulnerability Operations Center}}}
\fancyhead[R]{\small\textcolor{vocdark}{Rapport de stage}}
\fancyfoot[C]{\small\thepage}
\renewcommand{\headrulewidth}{0.4pt}
\renewcommand{\headrule}{\hbox to\headwidth{\color{vocblue}\leaders\hrule height \headrulewidth\hfill}}

% ---------------------------------------------------------------------------
% SECTION FORMATTING
% ---------------------------------------------------------------------------
\titleformat{\section}
  {\Large\bfseries\color{vocdark}}
  {\thesection}{0.8em}{}
\titleformat{\subsection}
  {\large\bfseries\color{vocblue}}
  {\thesubsection}{0.7em}{}
\titleformat{\subsubsection}
  {\normalsize\bfseries\color{vocdark}}
  {\thesubsubsection}{0.6em}{}

\captionsetup{font=small,labelfont={bf,color=vocblue}}

% ---------------------------------------------------------------------------
% CODE STYLE
% ---------------------------------------------------------------------------
\lstset{
  basicstyle=\ttfamily\small,
  keywordstyle=\color{vocblue}\bfseries,
  commentstyle=\color{gray},
  stringstyle=\color{vocgreen},
  numbers=left,
  numberstyle=\tiny\color{gray},
  frame=single,
  rulecolor=\color{voclight},
  breaklines=true,
  backgroundcolor=\color{white},
  tabsize=2
}

% ---------------------------------------------------------------------------
% COMMANDS
% ---------------------------------------------------------------------------
\newcommand{\tool}[1]{\textbf{\texttt{#1}}}
\newcommand{\techicon}[1]{\raisebox{-0.12em}{\includegraphics[height=1.0em]{icons/#1}}}
\newcommand{\techblock}[4]{%
  \par\medskip
  \noindent
  \begin{minipage}[t]{0.065\linewidth}\raggedleft\includegraphics[height=1.7em]{icons/#1}\end{minipage}%
  \hspace{0.45em}%
  \begin{minipage}[t]{0.88\linewidth}
    {\large\bfseries\color{vocblue}#2}\\[2pt]
    {\small\textbf{Description :} #3}\\
    {\small\textbf{Pourquoi ce choix :} #4}
  \end{minipage}
  \par
}

\begin{document}

% ===========================================================================
% PAGE DE GARDE
% ===========================================================================
\begin{titlepage}
\centering

{\LARGE\bfseries\color{vocdark} TEK-UP University}\\[0.25cm]
{\large \'Ecole d'Ing\'enierie \& Universit\'e priv\'ee}\\[0.2cm]
{\color{vocblue}\rule{\linewidth}{1.2pt}}\\[1.2cm]

{\Large\bfseries Soci\'et\'e : G\'erence Informatique}\\[0.5cm]

\vfill

{\large\color{vocdark} Rapport de stage}\\[0.3cm]
{\Huge\bfseries\color{vocdark} VOC Platform}\\[0.5cm]
{\LARGE\color{vocblue} Vulnerability Operations Center}\\[0.6cm]
{\large Plateforme automatis\'ee de gestion des vuln\'erabilit\'es}

\vfill

\begin{minipage}{0.8\linewidth}
\centering\small
\begin{tabular}{rl}
\toprule
\textbf{R\'ealis\'e par} & TBINI Mustapha Amin \\[4pt]
\textbf{Encadr\'ement} & Mr. MDAOUKHI Adel \\[4pt]
\textbf{Soci\'et\'e d'accueil} & G\'erence Informatique \\[4pt]
\textbf{Universit\'e} & TEK-UP University \\[4pt]
\textbf{Ann\'ee universitaire} & 2025--2026 \\[4pt]
\textbf{Lieu de d\'eploiement} & \texttt{192.168.184.135} (VMware) \\[4pt]
\textbf{P\'eriode} & Stage universitaire -- ann\'ee 2025--2026 \\
\bottomrule
\end{tabular}
\end{minipage}

\vfill
{\small\textcolor{gray} \today}
\end{titlepage}

% ===========================================================================
% REMERCIEMENTS
% ===========================================================================
\section*{Remerciements}
\addcontentsline{toc}{section}{Remerciements}

Au terme de ce stage, je tiens \`a exprimer ma gratitude \`a toutes les personnes
qui ont contribu\'e \`a la r\'ealisation de ce projet.

Je remercie tout d'abord mon encadrant, \textbf{Mr.~MDAOUKHI Adel}, pour son
accompagnement, ses conseils techniques et la confiance qu'il m'a accord\'ee tout
au long de ce travail.

Je remercie \'egalement la soci\'et\'e \textbf{G\'erence Informatique} pour son
accueil, ainsi que l'ensemble de ses \'equipes pour leur disponibilit\'e et leur
entraide.

Enfin, je remercie l'universit\'e \textbf{TEK-UP University} et son corps
enseignant pour la qualit\'e de la formation dispens\'ee durant l'ann\'ee
universitaire 2025--2026.

\clearpage

% ===========================================================================
% TABLE DES MATIERES
% ===========================================================================
\tableofcontents
\clearpage

% ===========================================================================
% 1. INTRODUCTION GENERALE
% ===========================================================================
\section{Introduction g\'en\'erale}

\subsection{Contexte}

La s\'ecurit\'e des syst\`emes d'information est devenue un enjeu
strat\'egique pour toutes les organisations. Chaque ann\'ee, des dizaines de
milliers de vuln\'erabilit\'es sont publi\'ees dans les bases de r\'ef\'erence
(CVE), et un nombre croissant d'entre elles sont activement exploit\'ees avant
m\^eme la disponibilit\'e d'un correctif. Le catalogue CISA KEV recense les
vuln\'erabilit\'es d\'ej\`a exploit\'ees dans la nature, et les outils
d'analyse estiment que seule une fraction des CVE constitue un risque r\'eel
pour une organisation donn\'ee.

Face \`a ce constat, la simple collecte de vuln\'erabilit\'es ne suffit plus.
Les \'equipes de s\'ecurit\'e doivent \^etre capables de :

\begin{itemize}[leftmargin=1.6em]
  \item \textbf{prioriser} efficacement les failles \`a traiter, en fonction
        du risque r\'eel pour l'activit\'e ;
  \item \textbf{corr\'eler} les d\'ecouvertes avec le contexte des menaces
        (exploits publics, exploitation active, indicateurs de compromission) ;
  \item \textbf{orchestrer} le traitement : ouverture de tickets, suivi de la
        rem\'ediation, mesure de l'am\'elioration.
\end{itemize}

\subsection{Pr\'esentation du projet}

Le projet d\'evelopp\'e durant ce stage, la \textbf{VOC Platform}
(\emph{Vulnerability Operations Center}), est une plateforme automatis\'ee de
gestion des vuln\'erabilit\'es qui :

\begin{itemize}[leftmargin=1.6em]
  \item scanne p\'eriodiquement le r\'eseau interne (\textbf{tous les ports}
        TCP, avec d\'etection des services et des syst\`emes d'exploitation) ;
  \item enrichit chaque d\'ecouverte avec de l'\textbf{intelligence sur les
        menaces} (NVD, MISP, EPSS, CISA KEV, OSV, Exploit-DB) ;
  \item calcule un \textbf{score de risque} orient\'e m\'etier entre 0 et 10 ;
  \item publie les r\'esultats dans une \textbf{stack de visualisation}
        (Elasticsearch + Kibana) ;
  \item ouvre automatiquement des \textbf{tickets de rem\'ediation} dans un
        outil ITSM (GLPI) ;
  \item partage les nouvelles d\'ecouvertes sous forme d'\textbf{\'ev\'enements
        MISP} ;
  \item supervise l'ensemble via une \textbf{observabilit\'e unifi\'ee}
        (Fleet + Beats).
\end{itemize}

\subsection{Organisation du rapport}

Ce rapport est organis\'e comme suit :

\begin{enumerate}[leftmargin=1.6em]
  \item la \textbf{probl\'ematique} \`a laquelle la plateforme apporte une
        solution, puis la pr\'esentation de l'\textbf{environnement de stage} ;
  \item la pr\'esentation des \textbf{outils et technologies} employ\'es et
        les raisons de leur choix ;
  \item l'\textbf{architecture} de la plateforme et sa \textbf{mod\'elisation
        UML} (cas d'utilisation, classes, s\'equence, \'etats, d\'eploiement) ;
  \item le \textbf{pipeline} de traitement des vuln\'erabilit\'es et le
        \textbf{moteur de calcul du risque} ;
  \item le \textbf{mod\`ele de donn\'ees}, la \textbf{s\'ecurit\'e},
        l'\textbf{\'etat op\'erationnel}, les \textbf{difficult\'es}
        rencontr\'ees et les \textbf{tests} ;
  \item enfin, la \textbf{conclusion} g\'en\'erale du stage.
\end{enumerate}
